\section{Quantitative intrinsic nondegeneracy and finite-data stability}
\label{sec:v24-quantitative-gluing}

Exact analytic continuation and exact signature matching are rigidity
mechanisms; by themselves they are not stability statements.  We isolate here
the quantitative hypotheses used in the statistical global theorem and derive
a modulus which is uniform on the resulting analytic class.  This separates
the exact inverse theorem from the finite-sample stability statement.

Fix a marked lifted-channel graph together with a designated rank-two
anchoring pair and a rooted signature-rigid spanning tree.  Let
\(\mathfrak K\) be a family of labelled periodic tables modulo one global
\(\mathrm{SE}(2)\) motion.  We say that \(\mathfrak K\) is a
\emph{uniformly nondegenerate analytic class} if the following constants can
be chosen independently of \(T\in\mathfrak K\).
\begin{enumerate}
\item There are \(\rho,B>0\) such that every oriented arclength curvature
function occurring in a designated obstacle chart extends holomorphically to
\(\{s+it:|t|<\rho\}\) and has absolute value at most \(B\) there.  The real
boundaries have uniform positive curvature, separation, diameter and channel
collar bounds.
\item There is a channel-persistence margin \(\chi>0\): every designated
closest-pair contact remains in the interior of its chosen patches, and the
length of every competing contact pair allowed by the finite labelled graph
exceeds the designated channel length by at least \(\chi\).
\item For each incidence used by the anchoring cycles or spanning tree there
are an integer \(m_{\rm sig}\) and a number \(\sigma_{\rm sig}>0\) such that
the finite signature
\[
 \mathfrak j^{[m_{\rm sig}]}(p)
   =(\kappa(p),\partial_s\kappa(p),\ldots,
                         \partial_s^{m_{\rm sig}}\kappa(p))
\]
is at distance at least \(\sigma_{\rm sig}\) from the corresponding finite
signature at every competing framed point outside the prescribed matching
arc, after quotienting only by the explicitly declared exact Euclidean
symmetries of that obstacle.
\item If \(v_1(T),v_2(T)\) are the two data-determined holonomy vectors of the
rank-two anchoring pair, then
\begin{equation}\label{eq:v24-holonomy-margin}
 \sigma_{\min}\bigl(v_1(T)\ \ v_2(T)\bigr)
                       \ge\sigma_{\rm hol}>0.
\end{equation}
\end{enumerate}
We equip the class with the labelled \(C^q\) metrics \(d_q\) after a fixed
root-frame gauge.  The holomorphic-collar hypothesis makes the corresponding
quotient class precompact in every fixed \(C^q\) topology; throughout this
section \(\mathfrak K\) is assumed closed and hence compact.

For an edge \(e\), normalize the action derivatives by
\[
 x_{e,b,m}(T)=\frac{\rho^m}{m!}S_{e,b}^{(m)}(0;T),
 \qquad m\ge2,
\]
and include the gaps after dividing by a fixed positive length scale
\(g_*\).  Let \(\mathbf x(T)\) be the resulting countable coordinate vector,
and set
\begin{equation}\label{eq:v24-data-metric}
 d_{\mathcal D}(T,T')
 =\sum_{e,b}\sum_{m\ge2}2^{-(m+|e|+3)}
       \min\{1,|x_{e,b,m}(T)-x_{e,b,m}(T')|\}
  +\sum_e2^{-(|e|+2)}
       \min\{1,|g_e(T)-g_e(T')|/g_*\}.
\end{equation}
Any fixed enumeration \(|e|\) of the finite measured edge set may be used.
Cauchy estimates on the fixed holomorphic collar make this metric compatible
with coordinatewise convergence of the complete intrinsic action data.

\begin{lemma}[Persistence of the prescribed gluing]
\label{lem:v24-gluing-persistence}
On a uniformly nondegenerate analytic class, the incidence matches selected by
the anchoring pair and spanning tree are locally unique and depend
continuously on the recovered finite signatures.  The recovered lattice map
\(L\) depends locally Lipschitz-continuously on the two holonomy vectors, with
a Lipschitz constant bounded in terms of \(\sigma_{\rm hol}\) and the marked
deck pair.  In particular, channel switching, signature switching and loss of
rank-two holonomy cannot occur inside a sufficiently small neighborhood whose
size is uniform on \(\mathfrak K\).
\end{lemma}

\begin{proof}
The signature gap \(\sigma_{\rm sig}\) makes the designated finite-signature
match the unique minimizer among competing framed points; continuity of the
finite jet under the holomorphic-collar topology preserves a gap of at least
\(\sigma_{\rm sig}/2\) under a uniform small perturbation.  The channel margin
\(\chi\) gives the analogous persistence of the selected closest pair.
For the lattice, write \(L=VM^{-1}\), where
\(V=(v_1\ v_2)\) and the fixed invertible matrix \(M\) has columns the two
marked deck vectors.  Matrix inversion is not being applied to noisy \(V\);
only multiplication by the fixed \(M^{-1}\) is needed.  The lower singular
value in \eqref{eq:v24-holonomy-margin} keeps the realized lattice uniformly
nondegenerate and gives uniform control of all subsequent deck corrections.
The finite graph then propagates the same matching choices along the rooted
tree.
\end{proof}

\begin{theorem}[Uniform inverse modulus on the robust analytic class]
\label{thm:v24-global-modulus}
For every integer \(q\ge0\) define
\begin{equation}\label{eq:v24-global-modulus-def}
 \omega_q(t)=
 \sup\{d_q(T,T'):T,T'\in\mathfrak K,
                         d_{\mathcal D}(T,T')\le t\}.
\end{equation}
Then \(\omega_q(t)<\infty\), it is nondecreasing, and
\begin{equation}\label{eq:v24-global-modulus-zero}
                         \omega_q(t)\longrightarrow0
                         \qquad(t\downarrow0).
\end{equation}
Moreover there is a deterministic sequence \(r_M\downarrow0\), depending
only on the fixed holomorphic collar and the finite measured graph, such that
if \(\mathcal D_M\) denotes the gaps and action jets through order \(M\), then
\begin{equation}\label{eq:v24-finite-modulus}
 d_q(T,T')
 \le
 \omega_q\!\left(
   C\|\mathcal D_M(T)-\mathcal D_M(T')\|_{M}+r_M
              \right),
 \qquad T,T'\in\mathfrak K.
\end{equation}
Here \(\|\cdot\|_M\) is the corresponding normalized finite-coordinate norm.
One may take \(r_M=C2^{-M}\) after increasing \(C\).

Thus the finite-coordinate resolution of the analytic inverse is accompanied
by an actual uniform modulus on the declared robust class.  No claim of a
universal modulus outside a fixed holomorphic collar is made.
\end{theorem}

\begin{proof}
The complete intrinsic data are injective on \(\mathfrak K\): the all-order
signed inverse recovers every incident analytic boundary, the rank-two
holonomy theorem recovers the lattice metric and realization, and
Lemma~\ref{lem:v24-gluing-persistence} selects the prescribed finite gluing.
The map \(T\mapsto\mathbf x(T)\) is continuous.  A continuous injective map
from the compact space \(\mathfrak K\) into the metric data space is a
homeomorphism onto its image.  Its inverse is therefore uniformly continuous,
which is exactly \eqref{eq:v24-global-modulus-zero}.

It remains to compare the full and truncated data metrics.  The terms with
\(m>M\) in \eqref{eq:v24-data-metric} contribute at most a constant times
\(2^{-M}\), independently of the values of those coordinates.  The finite
part is bounded by a constant times the normalized finite-coordinate norm.
Hence
\[
 d_{\mathcal D}(T,T')
 \le C\|\mathcal D_M(T)-\mathcal D_M(T')\|_M+C2^{-M}.
\]
Insert this estimate into \eqref{eq:v24-global-modulus-def}.
\end{proof}

\begin{remark}[Exact continuation versus stable reconstruction]
\label{rem:v24-exact-versus-stable}
The identity theorem is used only for exact-data uniqueness.  Statistical
stability enters through the explicitly compact holomorphic-collar class and
the modulus \(\omega_q\); the argument does not assert that numerical analytic
continuation is well conditioned on unrestricted analytic functions.  The
margins above also make precise why near-repeated signatures or a nearly
singular holonomy pair are excluded from the uniform statistical theorem
rather than hidden inside a singleton-gluing assumption.
\end{remark}
